\documentclass[a4paper,11pt]{article}

\input{config}
% ==================================================
% Document
% ==================================================

\begin{document}

% ==================================================
% Title Page
% ==================================================

\begin{titlepage}

\centering

\vspace*{1.5in}

{\Large ESE 507\par}

\vspace{0.5in}

{\Huge\bfseries Project: Part 1\par}

\vspace{1in}

{\large
Jin Yuan Chen\\[0.25em]
115702978
\par}

\vfill

{\large \today\par}

\end{titlepage}

% \section{Introduction}

% \subsection{Overview}

% \subsection{Specification}

% \subsection{Architecture}


\tableofcontents

\newpage

\section{Multiply Accumulate (MAC) Design}

Two versions of the MAC were implemented to compare the unpipelined and pipelined designs. Both the unpipelined \texttt{mac} and pipelined \texttt{mac\_pipe} designs use an accumulator register that updates its output on the rising edge of the clock. The primary difference between the two designs is the combinational multiply-add operation. In the unpipelined \texttt{mac}, the multiplication and addition are performed within a single combinational stage before the result is stored in the accumulator. In the pipelined \texttt{mac\_pipe}, the multiply-add operation is divided into two stages. The multiplier output and associated control signals are stored in an additional intermediate register before the addition is performed and the result is stored in the accumulator on the following clock cycle. Consequently, the pipelined design introduces an additional cycle of latency, requiring two clock cycles for an input to propagate to the accumulator output, compared with one cycle for the unpipelined design. However, the benefits of pipelining become apparent when considering the overall throughput and operating frequency of the system. 

\subsection{Design Parameter}
Both the \texttt{mac} and \texttt{mac\_pipe} designs were synthesized with \texttt{WIDTH = 16} and \texttt{ACCW = 48} across a range of clock frequencies. The analysis spans frequencies from a relatively low frequency of 10~MHz to a maximum of 1~GHz. At each frequency, the area, power, critical path location, and timing constraint of each design are evaluated.

\newpage
\section{Combinational Synthesis Results}

The unpipelined \texttt{mac} design was synthesized to determine the maximum achievable clock frequency with \texttt{WIDTH = 16} and \texttt{ACCW = 48}. The synthesis workflow was automated with scripts that write the configured design parameters and the synthesis script, before running synthesis on the remote CAD server. The generated synthesis logs are filtered for key attributes, including area, power, and timing. For convenient data analysis, these attributes are stored in comma-separated value (CSV) files for every synthesis run. The tabular structure of the CSV files is well suited for Python plotting libraries, which were used to generate the graphs below. For more details on the automation workflow, refer to \texttt{scr/}. Only data points that \textbf{meet the timing constraint} are plotted in the graphs below.

% For each frequency you try, record the area, power, the critical path location, and whether the timing 
% constraint was met or violated.

% make a table that shows this data for each attempted 
% frequency.  

%Make  sure  you  include  units  on  all  values  you  report 

\subsection{Area}

The design area estimated from the synthesis log is reported as the total cell area and does not account for net interconnect area. Therefore, the actual physical area of the design may vary depending on the interconnect area introduced during physical implementation. For the purposes of this analysis, the total design area is defined as the total cell area reported in the synthesis log.

As shown in Figure~\ref{fig:frequency-area}, the total design area remains relatively constant as the operating frequency increases. However, an overall upward trend can be observed, particularly at frequencies above 200~MHz. This increase can be attributed to the additional area required to meet the more difficult timing constraints associated with higher operating frequencies. The approximately linear increase in area beyond 200~MHz indicates the achievable trade-off for frequency at the expensive of area. In contrast, the relatively constant area at lower frequencies suggests that the synthesis tool implemented the design to a near-minimal cell configuration regardless of less restrictive timing constraints. 

\begin{figure}[!htbp]
    \centering

    \begin{multicols}{2}

        % Figure
        \centering
        \includesvg[inkscapelatex=false,width=\linewidth]{figures/mac_freq_area.svg}

        \captionof{figure}{Design area with increasing frequency.}
        \label{fig:frequency-area}

        \columnbreak

        % Table
        \centering
        \begin{tabular}{rr}
            \toprule
            \textbf{Frequency (MHz)} & \textbf{Area ($\mu m^2$)} \\
            \midrule
            10     & 2067.09 \\
            20     & 2067.09 \\
            25     & 2067.09 \\
            50     & 2067.09 \\
            100    & 2067.09 \\
            166.67 & 2067.09 \\
            200    & 2066.29 \\
            $\vdots$ & $\vdots$ \\
            500    & 2320.85 \\
            526.32 & 3200.90 \\
            555.56 & 2302.76 \\
            769.23 & 2424.06 \\
            \bottomrule
        \end{tabular}

        \label{tab:frequency-area}

    \end{multicols}
\end{figure}
 

An interesting pattern is the occasional small downward trend in area at higher frequencies. This suggests that, at certain frequency, the synthesis tool may identify a simplified version that satisfies the required timing constrain with less cell area. For example, the design synthesized at 555.56~MHz has a slightly lower total area than the design synthesized at 500~MHz, despite the overall increasing trend between operating frequency and area. This behavior indicates that the relationship between frequency and synthesized area is not strictly monotonic and can depend on the specific optimizations selected by the synthesis tool.

\subsection{Power}

The Total Power shown in Figure 2 is defined as the sum of Total Dynamic Power and Cell Leakage Power. The Total Dynamic Power reported in the synthesis log is also the sum of Cell Internal Power and Net Switching Power, both contributing an near equal portion of the sum. 

\begin{figure}[htbp]
    \centering

    \begin{multicols}{2}

        \centering
        \includesvg[inkscapelatex=false,width=\linewidth]{figures/mac_freq_power.svg}
        \captionof{figure}{Dynamic power with increasing frequency.}
        \label{fig:dynamic-power}

        \columnbreak

        \centering

        \begin{tabular}{rr}
            \toprule
            \textbf{Frequency (MHz)} & \textbf{$\sum$ Power ($\mu W$)} \\
            \midrule
            10     & 65.49 \\
            20     & 83.85 \\
            25     & 93.03 \\
            50     & 138.92 \\
            100    & 230.72 \\
            166.67 & 353.12 \\
            200    & 411.86 \\
            $\vdots$ & $\vdots$ \\
            500    & 1127.93 \\
            526.32 & 1148.45 \\
            555.56 & 1218.61 \\
            769.23 & 1745.19 \\
            \bottomrule
        \end{tabular}
    \end{multicols}
\end{figure}

The Cell Leakage Power is almost negligible compared with the dynamic power. This may be attributed to the relatively small area of the design, with the accumulator register being the primary memory elements contributing to leakage power. In the case that the operating frequency is considerably slow, the dynamic power can become sufficiently small that leakage power becomes a significant contributor to total power. Otherwise, given a fixed accumulator width and integer-bit width, the leakage power remains approximately constant at frequency below 166.67 MHz, as shown in the table of Figure 1. This suggests that leakage power is primarily dependent on the area of the design. The slight increase in leakage power at higher frequencies can therefore be attributed to the corresponding increase in design area.

Consequently, the trend observed in Figure 2 for total power is largely influenced by the Dynamic Power and its dependence on operating frequency. Increasing the frequency results in a near linear increase in Dynamic Power, with an approximate power increase of 1.836 $\mu\text{W}/\text{MHz}$. In other words, a higher-frequency design performs more switching operation per unit time than a lower-frequency design, resulting in higher power consumption over the same period of time. 


\subsection{Critical Paths}

The synthesis report provides the start and end points of the critical paths, along with the paths themselves and their data arrival times. If the data arrival time is less than or equal to the data required time, the design constraint is considered MET, as indicated by Slack (MET) in the report.

A design with Slack (VIOLATED) indicates a timing constraint violation and will typically display a negative slack value, calculated as the difference between the data arrival time and the data required time. A negative value indicates that the timing constraint has not been met. A design reported as Slack (VIOLATED: increase significant digits) is also considered a violation, even if the displayed slack value is zero. This occurs because the actual slack value is slightly negative but is rounded to zero due to the limited number of significant digits displayed in the report.

\begin{table}[htbp]
    \centering
    \scriptsize
    \begin{tabular}{r r r l l l}
        \toprule
        \textbf{Frequency (MHz)} &
        \textbf{Area ($\mu m^2$)} &
        \textbf{Slack (ns)} &
        \textbf{Path Start} &
        \textbf{Path End} &
        \textbf{Timing} \\
        \midrule
        10      & 2067.09 & 95.68  & \texttt{input0[1]}  & \texttt{acc\_reg[47]} & MET \\
        11.11   & 2067.09 & 85.68  & \texttt{input0[1]}  & \texttt{acc\_reg[47]} & MET \\
        12.5    & 2067.09 & 75.68  & \texttt{input0[1]}  & \texttt{acc\_reg[47]} & MET \\
        14.29   & 2067.09 & 65.68  & \texttt{input0[1]}  & \texttt{acc\_reg[47]} & MET \\
        16.67   & 2067.09 & 55.68  & \texttt{input0[1]}  & \texttt{acc\_reg[47]} & MET \\
        20      & 2067.09 & 45.68  & \texttt{input0[1]}  & \texttt{acc\_reg[47]} & MET \\
        25      & 2067.09 & 35.68  & \texttt{input0[1]}  & \texttt{acc\_reg[47]} & MET \\
        33.33   & 2067.09 & 25.68  & \texttt{input0[1]}  & \texttt{acc\_reg[47]} & MET \\
        50      & 2067.09 & 15.68  & \texttt{input0[1]}  & \texttt{acc\_reg[47]} & MET \\
        100     & 2067.09 & 5.68   & \texttt{input0[1]}  & \texttt{acc\_reg[47]} & MET \\
        111.11  & 2067.09 & 4.68   & \texttt{input0[1]}  & \texttt{acc\_reg[47]} & MET \\
        125     & 2067.09 & 3.68   & \texttt{input0[1]}  & \texttt{acc\_reg[47]} & MET \\
        142.86  & 2067.09 & 2.68   & \texttt{input0[1]}  & \texttt{acc\_reg[47]} & MET \\
        166.67  & 2067.09 & 1.68   & \texttt{input0[1]}  & \texttt{acc\_reg[47]} & MET \\
        200     & 2066.29 & 0.68   & \texttt{input0[1]}  & \texttt{acc\_reg[47]} & MET \\
        250     & 2102.20 & 0.01   & \texttt{input0[0]}  & \texttt{acc\_reg[47]} & MET \\
        333.33  & 2164.18 & 0.00   & \texttt{input0[7]}  & \texttt{acc\_reg[47]} & MET \\
        500     & 2320.85 & 0.00   & \texttt{input0[3]}  & \texttt{acc\_reg[47]} & MET \\
        526.32  & 2300.90 & 0.03   & \texttt{input0[5]}  & \texttt{acc\_reg[47]} & MET \\
        555.56  & 2302.76 & 0.00   & \texttt{input0[7]}  & \texttt{acc\_reg[47]} & MET \\
        588.24  & 2295.58 & 0.01   & \texttt{input0[11]} & \texttt{acc\_reg[47]} & MET \\
        625     & 2321.12 & 0.00   & \texttt{input0[5]}  & \texttt{acc\_reg[47]} & MET \\
        666.67  & 2341.60 & 0.00   & \texttt{input0[5]}  & \texttt{acc\_reg[44]} & MET \\
        714.29  & 2395.60 & 0.00   & \texttt{input0[0]}  & \texttt{acc\_reg[35]} & MET \\
        769.23  & 2424.06 & 0.00   & \texttt{input0[7]}  & \texttt{acc\_reg[31]} & MET \\
        833.33  & 2471.14 & $-0.07$ & \texttt{input0[5]}  & \texttt{acc\_reg[33]} & VIOLATED \\
        909.09  & 2543.23 & $-0.14$ & \texttt{input0[15]} & \texttt{acc\_reg[21]} & VIOLATED \\
        1000    & 2529.13 & $-0.22$ & \texttt{input0[7]}  & \texttt{acc\_reg[38]} & VIOLATED \\
        \bottomrule
    \end{tabular}
    \caption{Frequency, area, timing slack, and critical path for unpipelined \texttt{mac}.}
    \label{tab:frequency-area-slack-path}
\end{table}

The critical path starts from an \texttt{input0} register and ends at \texttt{acc\_reg}, with slight variations in the indices of \texttt{input0} and \texttt{acc\_reg}. At low frequencies and up to approximately 200~MHz, the critical path remains the same, starting from \texttt{input0[1]} and ending at \texttt{acc\_reg[47]}. Between 200~MHz and 625~MHz, the endpoint of the critical path remains fixed at \texttt{acc\_reg[47]}, while the index of \texttt{input0} varies depending on the operating frequency. Above 625~MHz, both the \texttt{input0} index and the endpoint index of \texttt{acc\_reg} vary without a clear or consistent pattern as shown in Table 1.

In particular, there appears to be a tendency for the critical path to shift toward higher-indexed \texttt{input0} elements and, in other cases, lower-indexed \texttt{acc\_reg} elements. This indicate a possible trend in which the synthesis process optimizes the design by shifting the critical path toward higher-indexed \texttt{input0} elements. However, additional analysis is required to determine whether this trend is directly caused by the stricter timing constraints, the design area, or their interaction. 

%Q2
\section{Pipelined Synthesis Results}

In a similar fashion, the pipelined \textit{mac\_pipe} design was synthesized to determine the maximum achieavable clock frequency with WIDTH = 16 and ACCW = 48. The same automated scripts were used to collect and store the synthesis results in a CSV file. The resulting datasets for the unpipelined \textit{mac} and pipelined \textit{mac\_pipe} designs were then analyzed to compare their area, power, and timing performance. In terms of performance metric, the maximum achievable operating frequency was prioritized before considering tradeoff in area and power, whose relative importance depends on the target application. Only data points
that meet the timing constraint are plotted in the graphs below.


% For each frequency you try, record the area, power, the critical path location, and whether the timing 
% constraint was met or violated.

% make a table that shows this data for each attempted 
% frequency.  

%Make  sure  you  include  units  on  all  values  you  report 

\subsection{Area}

\begin{figure}[htbp]
    \centering

    \begin{multicols}{2}

        \centering
        \includesvg[inkscapelatex=false,width=\linewidth]{figures/mac_pipe_freq_area.svg}
        \captionof{figure}{Dynamic power with increasing frequency.}
        \label{fig:pipe_area}

        \columnbreak

        \centering
        \begin{tabular}{rrr}
            \toprule
            \textbf{Freq.\ (MHz)} &
            \multicolumn{2}{c}{\textbf{Area ($\mu m^2$)}} \\
            \cmidrule(lr){2-3}
            & \textbf{\texttt{mac}} & \textbf{\texttt{mac\_pipe}} \\
            \midrule
            10     & 2067.09 & 2306.75 \\
            20     & 2067.09 & 2306.49 \\
            25     & 2067.09 & 2306.49 \\
            50     & 2067.09 & 2306.49 \\
            100    & 2067.09 & 2306.49 \\
            166.67 & 2067.09 & 2306.49 \\
            200    & 2066.29 & 2305.16 \\
            $\vdots$ & $\vdots$ & $\vdots$ \\
            500    & 2320.85 & 2648.03 \\
            526.32 & 2300.90 & 2639.78 \\
            555.56 & 2302.76 & 2692.98 \\
            769.23 & 2424.06 & 2773.58 \\
            \bottomrule
        \end{tabular}
    \end{multicols}
\end{figure}


% Make two graphs. On both, show 
% clock frequency on the x­axis; then show area as the y­axis on one graph and power as the y­axis on 
% the other.)

% Explain the trends that you observed and explain why they occur. (

% Only include the design points where the timing constraint is MET

\section{Power}
\begin{figure}[htbp]
    \centering

    \begin{multicols}{2}

        \centering
        \includesvg[inkscapelatex=false,width=\linewidth]{figures/mac_pipe_freq_power.svg}
        \captionof{figure}{Dynamic power with increasing frequency.}
        \label{fig:pipe_power}

        \columnbreak

        \centering

        \begin{tabular}{rrr}
            \toprule
            \textbf{Freq.\ (MHz)} &
            \multicolumn{2}{c}{\textbf{$\sum$ Power ($\mu W$)}} \\
            \cmidrule(lr){2-3}
            & \textbf{\texttt{mac}} & \textbf{\texttt{mac\_pipe}} \\
            \midrule
            10     & 65.49 & 71.03 \\
            20     & 83.85 & 90.75 \\
            25     & 93.03 & 100.61 \\
            50     & 138.92 & 149.92 \\
            100    & 230.72 & 248.53 \\
            166.67 & 353.12 & 380.01 \\
            200    & 411.86 & 446.49 \\
            $\vdots$ & $\vdots$ & $\vdots$ \\
            500    & 1127.93 & 1182.18 \\
            526.32 & 1148.45 & 1239.58 \\
            555.56 & 1218.61 & 1333.95 \\
            769.23 & 1745.19 & 1928.30 \\
            \bottomrule
        \end{tabular}
    \end{multicols}
\end{figure}

% Make two graphs. On both, show 
% clock frequency on the x­axis; then show area as the y­axis on one graph and power as the y­axis on 
% the other.)

% Explain the trends that you observed and explain why they occur. (

% Only include the design points where the timing constraint is MET

\subsection{Critical Paths}
% For each frequency, give a description in your report of where the critical path is. explain logically where the critical path lies in the module.

\begin{table}[htbp]
    \centering
    \scriptsize
    \begin{tabular}{r r r l l l}
        \toprule
        \textbf{Frequency (MHz)} &
        \textbf{Area ($\mu m^2$)} &
        \textbf{Slack (ns)} &
        \textbf{Path Start} &
        \textbf{Path End} &
        \textbf{Timing} \\
        \midrule
        10      & 2306.75 & 95.74  & \texttt{reg2/out\_accum\_reg[0]}  & \texttt{reg2/out\_accum\_reg[47]} & MET \\
        11.11   & 2306.75 & 85.74  & \texttt{reg2/out\_accum\_reg[0]}  & \texttt{reg2/out\_accum\_reg[47]} & MET \\
        12.5    & 2306.49 & 75.74  & \texttt{reg2/out\_accum\_reg[0]}  & \texttt{reg2/out\_accum\_reg[47]} & MET \\
        14.29   & 2306.75 & 65.74  & \texttt{reg2/out\_accum\_reg[0]}  & \texttt{reg2/out\_accum\_reg[47]} & MET \\
        16.67   & 2306.75 & 55.74  & \texttt{reg2/out\_accum\_reg[0]}  & \texttt{reg2/out\_accum\_reg[47]} & MET \\
        20      & 2306.49 & 45.74  & \texttt{reg2/out\_accum\_reg[0]}  & \texttt{reg2/out\_accum\_reg[47]} & MET \\
        25      & 2306.49 & 35.74  & \texttt{reg2/out\_accum\_reg[0]}  & \texttt{reg2/out\_accum\_reg[47]} & MET \\
        33.33   & 2306.49 & 25.74  & \texttt{reg2/out\_accum\_reg[0]}  & \texttt{reg2/out\_accum\_reg[47]} & MET \\
        50      & 2306.49 & 15.74  & \texttt{reg2/out\_accum\_reg[0]}  & \texttt{reg2/out\_accum\_reg[47]} & MET \\
        100     & 2306.49 & 5.74   & \texttt{reg2/out\_accum\_reg[0]}  & \texttt{reg2/out\_accum\_reg[47]} & MET \\
        111.11  & 2306.49 & 4.74   & \texttt{reg2/out\_accum\_reg[0]}  & \texttt{reg2/out\_accum\_reg[47]} & MET \\
        125     & 2306.49 & 3.74   & \texttt{reg2/out\_accum\_reg[0]}  & \texttt{reg2/out\_accum\_reg[47]} & MET \\
        142.86  & 2306.49 & 2.74   & \texttt{reg2/out\_accum\_reg[0]}  & \texttt{reg2/out\_accum\_reg[47]} & MET \\
        166.67  & 2306.49 & 1.74   & \texttt{reg2/out\_accum\_reg[0]}  & \texttt{reg2/out\_accum\_reg[47]} & MET \\
        200     & 2305.16 & 0.75   & \texttt{reg2/out\_accum\_reg[0]}  & \texttt{reg2/out\_accum\_reg[47]} & MET \\
        250     & 2342.93 & 0.12   & \texttt{reg1/out\_product\_reg[0]} & \texttt{reg2/out\_accum\_reg[47]} & MET \\
        333.33  & 2459.44 & 0.00   & \texttt{reg2/out\_accum\_reg[0]}  & \texttt{reg2/out\_accum\_reg[47]} & MET \\
        500     & 2648.03 & 0.00   & \texttt{reg2/out\_accum\_reg[9]}  & \texttt{reg2/out\_accum\_reg[47]} & MET \\
        526.32  & 2639.78 & 0.00   & \texttt{reg2/out\_accum\_reg[7]}  & \texttt{reg2/out\_accum\_reg[46]} & MET \\
        555.56  & 2692.98 & 0.00   & \texttt{input0[7]}  & \texttt{reg1/out\_product\_reg[30]} & MET \\
        588.24  & 2693.52 & 0.01   & \texttt{input0[3]}  & \texttt{reg1/out\_product\_reg[30]} & MET \\
        625     & 2722.78 & 0.00   & \texttt{input0[5]}  & \texttt{reg1/out\_product\_reg[30]} & MET \\
        666.67  & 2732.35 & 0.00   & \texttt{input0[5]}  & \texttt{reg1/out\_product\_reg[29]} & MET \\
        714.29  & 2724.64 & 0.00   & \texttt{input0[3]}  & \texttt{reg1/out\_product\_reg[30]} & MET \\
        769.23  & 2773.58 & 0.00   & \texttt{reg2/out\_accum\_reg[7]}  & \texttt{reg2/out\_accum\_reg[47]} & MET \\
        833.33  & 2823.86 & 0.00   & \texttt{input0[11]} & \texttt{reg1/out\_product\_reg[23]} & VIOLATED \\
        909.09  & 2860.56 & $-0.06$ & \texttt{input0[13]} & \texttt{reg1/out\_product\_reg[27]} & VIOLATED \\
        1000    & 2832.37 & $-0.18$ & \texttt{input0[9]}  & \texttt{reg1/out\_product\_reg[20]} & VIOLATED \\
        \bottomrule
    \end{tabular}
    \caption{Frequency, area, timing slack, and critical path for pipelined \texttt{mac\_pipe}.}
    \label{tab:pipe-frequency-area-slack-path}
\end{table}

%Q3
\section{Performance Comparison}
% Did pipelining help make this module faster? Explain why 
% or why not and show how this is reflected in the synthesis data and critical path

%Q4
\subsection{Maximum-Frequency Energy}

%  pipelined design with the maximum clock frequency you found, how much energy would your 
% system consume if it were to process a sequence of 50 sets of input values? 

% assume you have to wait 
% until the final output comes out of the system, and don’t forget that your pipelined design takes more 
% than 50 cycles to compute 50 sets of inputs.

% Remember: energy is measured in joules. Power = energy per time. 1 Watt = 1 Joule per second. Use 
% the power obtained from synthesis and your understanding of the time it would take for your system 
% to fully compute 50 sets of input values

%Q5
\subsection{Frequency-Dependent Energy}
% Would  the  energy  you  computed  in  question  3  change  if  you  resynthesized  the  design  targeting different clock frequencies? Explain and justify your answer. 

%Think carefully about what changes when you change the target frequency and how those changes affect the power the system consumes.

\subsection{Latency and Throughput}
% Make a table that compares the power, area, latency, and throughput of your pipelined and unpipelined 
% MAC designs

% In your report show how you calculated the latency and throughput. Quantify latency in seconds  (or  ns),  and  quantify  throughput  in  terms  of  MACs  per  second. TOPIC 6

\subsection{Design Trade-offs}

% Based  on  the  trade­offs  seen  in  your  table,  explain  when  it  would 
% make sense for a designer to choose the pipelined design and when it would make sense to use the 
% unpipelined design

%Q6
\section{Potential for Deeper Pipelining}

\subsection{Multiplier Pipelining}

% Based on your 
% results to questions 1 and 2, would you expect that deeper pipelining in the multiplier might help you 
% reach higher clock frequencies? Justify why or why not.

% If you were to pipeline the multiplier in this way, what other changes would you have to make in your 
% module? Be specific about each of the module’s control inputs: which of them would need to change, 
% and which would not?

\subsection{Adder Pipelining}
% Would pipelining inside of the adder be possible and a good idea? Why or why not?

% Answer this question based on your understanding of the design and your answers to prior questions. 
% You do not need to modify your design/code to answer this question.

%Q7
\section{Variable Parameter Synthesis Results}

% how changing WIDTH and ACCW affects the critical path location and maximum clock frequency of your pipelined MAC module

%Q7
\subsection{Accumulator Width (ACCW)}
% First, do a set of experiments where you set WIDTH=12 and synthesize three designs with ACCW=24, 48, 
% and 64

%Q7
\subsection{Numerical-Bit Width (WIDTH)}
% Next, do a new set of experiments where you set ACCW=48 and synthesize three designs with WIDTH=8, 
% 16, 24.

\subsection{Scaling Behavior}
% Do  the  clock  frequency  and  critical  path  location  change  as  ACCW  changes? 

% Do  the  clock  frequency  and  critical  path location change  when 
% WIDTH changes?

% Explain what you see and what you learn from it. You should expect to see differences in the scaling behavior of WIDTH and ACCW. 

% Explain why they behave differently. In your explanation, account for the shift­and­saturate logic on the output path as well as the multiplier and adder.

%Q8
\subsection{Accumulator Overflow}

%if the value stored in the accumulator grows large enough, it will overflow. That is, ACCW bits may not be enough to store the resulting number.

% Assume WIDTH=5 and ACCW=16. What is the maximum number of MACs your system could perform while guaranteeing that the accumulator cannot overflow?

% what is the largest magnitude number you could produce on the multiplier’s output? Then, how many accumulations would it take for that number to produce an overflow in the accumulator? 

% Don’t forget that our values are all signed integers, and  don’t  forget  that  the  accumulator  can  be  initialized  to  a  signed  WIDTH­bit  number.  Show  your reasoning and justify your answer.

\section{Appendix}
% Include the six synthesis reports for parameter scaling question along with your Final Report
















% \subsection{list the modules}

% \subsection{Functional Verification}

% \subsection{Logic Synthesis }




% \section{Memory input}

% \subsection{list the modules}

% \subsection{Functional Verification}

% \subsection{Logic Synthesis }

% \section{Control}

% \subsection{list the modules}

% \subsection{Functional Verification}

% \subsection{Logic Synthesis}

% \section{Speed Optimization}

% \subsection{Critical-Path Reduction}
% module logical implementation, unless extreme condition, will be optimized by synthesizer/complier

% \subsection{Pipeline Architecture}

% \subsection{Data-level Parallelism}

% \subsection{Instruction-Level Parallelism}

% \section{Discussion}

% \section{Conclusion}


\end{document}